\newpage
\section{Protokollmechanismen 2}
  \subsection{Einführung}
    Fehlerhäufigkeit
    \begin{itemize}
      \item Die Bitfehlerrate (bit error rate, BER) ist ein Maß für die Fehlerhäufigkeit. Sie wird
        durch das Verhältnis gestörter Bits zu übertragenen Bits berechnet:
        $$\text{Bitfehlerrate}=\frac{\sum\text{ gestörte Bits}}{\sum\text{ übertragene Bits}}$$
    \end{itemize}
    Redundanz
    \begin{itemize}
      \item Einführung von Redundanz um Fehler zu erkennen und ggfs zu beheben
        \begin{itemize}
          \item Fehlererkennung
          \begin{itemize}
            \item Redundanz zu Daten hinzufügen
            \item Verwende Codewörter, die sich hinreichend unterscheiden
            \item Beispiel
              \begin{itemize}
                \item "t" und "d" manchmal schwer unterscheidbar
                \item Verwende Nato-Alphabet "Tango" und "Delta"
              \end{itemize}
          \end{itemize}
        \item Fehlerkorrektur
          \begin{itemize}
            \item Falls Fehler erknannt, versuche mittels Redundanz zu korrigieren
          \end{itemize}
        \end{itemize}
    \end{itemize}
  \subsection{Erkennung von Bitfehlern}
    Erkennung von Bitfehlern
    \begin{itemize}
       \item Problem
         \begin{itemize}
           \item Wie können Bitfehler beim Empfänger erkannt werden?
         \end{itemize}
       \item Grundlegender Ansataz
         \begin{itemize}
           \item Hinzufügen redundanter Bits
             \begin{itemize}
               \item Bits, die nicht den Nutzdaten zugeordnet werden können und die nicht für
                 Kontrollzwecke benötigt werden
             \end{itemize}
         \end{itemize}
       \item Typische Beispiele in Rechnernetzen
         \begin{itemize}
           \item Paritätsbits
           \item Prüfsumme
         \end{itemize}
    \end{itemize}
    \subsubsection{Grundlagen}
      Hamming-Abstand (Hamming-Distanz)
      \begin{itemize}
        \item Es seien $C$ ein Code fester Länge und $v_1,v_2$ zwei Codewörter aus $C$. Der
          Hamming-Abstand $\Delta(v_1,v_2)$ ist die Anzahl an Positionen, an denen sich die
          Symbole in $v_1$ und $v_2$ unterschiedlich sind
          \begin{itemize}
            \item entspricht Anzahl der Eins-Bits von $v_1\oplus v_2$
          \end{itemize}
      \end{itemize}
      Code-Abstand
      \begin{itemize}
        \item Für einen Code $C$ fester Länge ist der Code-Abstand $\Delta C$ die minimale
          Distanz zwishen zwei Codewörtern
          $$\Delta C:=\min\{\Delta(v_1,v_2)\mid v_1,v_2\in C,v_1\neq v_2\}$$
      \end{itemize}
      Fehlererkennender, -korrigierender Code
      \begin{itemize}
        \item Ein Code ist $r$-fehlererkennend, wenn der Empfänger jedes Codewort, in dem
          bis zu $r$ Symbole verfälscht wurden, als fehlerhaft identifizert
        \item Ein Code ist $r$-fehlerkorrigierend, wenn der Empfänger jedes Codewort, in
          dem bis zu $r$ Symbole verfälscht wurden, wieder herstellten kann
      \end{itemize}
      Hamming-Kugel
      \begin{itemize}
        \item Es seien $C\subseteq\Pi^n$ ein Coded fester Länge, $v\in C$ und $r\in\nat$. Die Menge
          $$K_r(v):=\{w\in C\mid\Delta(v,w)\le r\}$$
          wird als Hamming-Kugel von Codewort $v$ mit Radius $r$ bezeichnet
        \item $K_r(v)$ enthält genau diejenigen Codewörter
          \begin{itemize}
            \item die sich von $v$ an höchstens $r$ Positionen unterscheiden
            \item d.h. deren Distanz zu $v$ beträgt höchstens $r$
          \end{itemize}
        \item Anschaulich
          \begin{itemize}
            \item Kugel um $v$ mit Radius $r$ im $n$-dimensionalen Raum
          \end{itemize}
      \end{itemize}
      Code ist nicht $r$-fehlererkennend
      \begin{itemize}
        \item Codewort einer Hamming-Kugel liebt
          \begin{itemize}
            \item auf dem Rand einer Hamming-Kugel eines anderen Codeworts
            \item oder innerhalb der Hamming-Kugel eines anderen Codewortes
            \item $\ra$ Codewort kann durch Ändern von $r$ oder weniger Bits in anderes, gültiges
              Codewort überführt werden
          \end{itemize}
      \end{itemize}
      Code ist $r$-fehlererkennend
      \begin{itemize}
        \item Code-ABstand größer als $r$
          \begin{itemize}
            \item d.h. Codewörter sind weiter als $r$ voneinander entfernt
            \item Hamming-Kugeln sind frei von anderen Codewörtern
            \item $\ra$ Codewort kann durch ändern von $r$ oder weniger Bits nicht in anderes,
              gültiges Codewort überführt werden
          \end{itemize}
      \end{itemize}
      Code ist nicht $r$-fehlerkorrigierend
      \begin{itemize}
        \item Zwei Hamming-Kugeln sind so nah zusammen, dass sie sich
          \begin{itemize}
            \item berühren
            \item oder schneiden
            \item $\ra$ Codewörter in der Schnittmenge können durch ändern von $r$ oder
              weniger Bits aus beiden Codewörtern entstanden sein
          \end{itemize}
      \end{itemize}
      Code ist $r$-fehlerkorrigierend
      \begin{itemize}
        \item Code-Abstand größer als $2r$
      \end{itemize}
      Fehlererkennd / -korrigierend
      \begin{itemize}
        \item Für einen Code $C$ fester Länge gilt
          \begin{itemize}
            \item $C$ ist $r$-fehlererkennend $\Leftrightarrow$ $\Delta C >r$
            \item $C$ ist $r$-fehlererkorrigierend $\Leftrightarrow$ $\Delta C >2r$
          \end{itemize}
      \end{itemize}
    \subsubsection{Paritätsbits}
      Einfache Paritätsprüfung
      \begin{itemize}
        \item Fehlererkennung von einzelnen Bitfehlern: Paritätsbit
          \begin{itemize}
            \item Nutzdaten: $x=(x_1,x_2,\dots,x_k)$
            \item Paritätsbits: $p=f(x)=(x_1+\dots+x_k)\mod 2$
            \item Codewort: $z=(xp)$
          \end{itemize}
      \end{itemize}
    \subsubsection{Hamming-Code}
      \begin{itemize}
        \item Eigenschaften
          \begin{itemize}
            \item Linearer $(n,k)$ Code
              \begin{itemize}
                \item $n$: Länge des Codewortes
                \item $k$: Anzahl der Nutzdatenbits im Codewort
                \item Anzahl der verwendeten Paritätsbits: $n-k$
              \end{itemize}
            \item Code-Abstand von 3
            \item Kann Einzelbitfehler korrigieren
            \item Mit einfachen Harwareschaltungen implementierbar
          \end{itemize}
      \end{itemize}
    \subsubsection{Cyclic Redundancy Check (CRC)}
      \begin{itemize}
        \item Ein linearer Code heißt zyklisch, wenn jede zyklische Verschiebung eines Codeworts
          ebenfalls ein Codewort ist
        \item Wenn $(a_1,a_2,\dots,a_n)$ Codewort von $C$ ist, dann sind dies auch\\
          $(a_2,a_3,\dots,a_n,a_1),(a_3,a_4,\dots,a_1,a_2),\dots$
        \item Zusammenhang zwischen Codewörtern und Polynomen
          \begin{itemize}
            \item Symbole eines Codeworts werden als Koeffizienten eines Polynoms interpretiert
          \end{itemize}
          \item Quellenaplphabet $\Sigma$ sei endlicher Körper $\mathbb{F}$
            \begin{itemize}
              \item Jedes Codewort $v$ der Länge $n$ lässt sich als Polynom $v(x)\in\mathbb{F}[x]$
                mit $\deg v(x)<n$ einem Codewort $v$ der Länge $n$
            \end{itemize}
          \item CRC häufig in der Praxis eingesetzt (z.B. Ethernet, WLAN)
      \end{itemize}
      Generatorpolynom
      \begin{itemize}
        \item Es seien $g(x)$ ein Polynom aus der Menge $\mathbb{F}[x]$ und $n\in\nat^+$. Der von
          $g(x)$ generierte Code der Länge $n$ ist die Menge
          $$C:=\{v(x)\mid\deg v/x)<n\text{ und }g(x)\text{ teilt }v(x)\text{ ohne Rest}\}$$
          $g(x)$ heißt das Generatorpolynom von $C$
        \item Das Generatorpolynom $g(x)$ erzeugt genau dann einen zyklischen Code der Länge $n$, wenn
          es die folgende Teilbarkeitsregel aufweist:
          $$g(x)\text{ teilt }x^n+1\text{ ohne Rest}$$
      \end{itemize}
      Berechnung CRC
      \begin{itemize}
        \item Wie werden die Codewörter erzeugt?
          \begin{itemize}
            \item Codewort der Länge $n$
            \item $r$ Prüfbits werden an die zu übertragenen Nutzdaten angehängt
              \begin{itemize}
                \item Damit $n-r$ Nutzdaten enthalten
              \end{itemize}
          \end{itemize}
        \item Seien
          \begin{itemize}
            \item $v(x)$: Codewortpolynom
            \item $u(x)$: Nutzdatenpolynom
            \item Dann gilt
              \begin{itemize}
                \item Die ersten $n-r$ Koeffizienten des Codewortpolynoms $v(x)$ müssen mit dem
                  Polynom $x^ru(x)$ übereinstimmen
              \end{itemize}
          \end{itemize}
      \end{itemize}
      CRC-Algorithmus (Sender)
      \begin{itemize}
        \item Gegeben: Generatorpolynom $g(x)$ mit $\deg g(x)=r$, Nutzdatenpolynom $u(x)$ mit 
          $\deg u(c)<n-r$
        \item Gesucht: Codewortpolynom $v(x)$ mit $\deg v(x)<n$
        \item Ablauf
          \begin{enumerate}[label=\arabic*)]
            \item Multipliziere Nutzdatenpolynom $u(x)$ mit $x^r$
            \item Dividiere $x^ru(x)$ durch Generatorpolynom $g(x)$
            \item Berechne Summe aus Divisionsrest $r(x)$ und $x^ru(x)$
          \end{enumerate}
      \end{itemize}
      CRC-Algorithmus (Empfänger)
      \begin{itemize}
        \item Gegeben: Generatorpolynom $g(x)$ mit $\deg g(x)=r$, Codewortpolynom $u(x)$ mit 
          $\deg u(c)<n-r$
        \item Gesucht: Nutzdatenpolynom $v(x)$ mit $\deg v(x)<n$
        \item Ablauf
          \begin{enumerate}[label=\arabic*)]
            \item Dividiere Empfangswort durch Generatorpolynom $g(x)$
            \item Kein Rest $\ra$ keine Fehler festgestellt\\
              Rest $\ra$ Übertragung war fehlerhaft
          \end{enumerate}
      \end{itemize}
      CRC-Implementierung in Hardware
      \begin{itemize}
        \item Realisierung in Hardware
          \begin{itemize}
            \item Benutzung von rückgekoppelten Schieberegistern
            \item CRC wird während des "Durchschiebens" durch das Schieberegister berechnet
          \end{itemize}
        \item Vorgehensweise
          \begin{itemize}
            \item Daten durchlaufen bitweise das Schieberegister
            \item XOR-Gatter an Stellen von "1"-Bits in $g(x)$
            \item An höchstem Bit Rückkopplung, falls "1" durchgeschoben wird
            \item Nach Durchschieben von $v(x)$ steht Prüfsumme im Register
              \begin{itemize}
                \item Zu beachten: Soll Prüfsumme ebenfalls "aus dem Register geschoben werden", ist
                  noch etwas mehr Schaltungslogik nötig
              \end{itemize}
          \end{itemize}
      \end{itemize}
  \subsection{Vorwärtsfehlerkorrektur}
    Vorwärtsfehlerkorrektur (Forward Error Correction -- FEC)
    \begin{itemize}
      \item Ziel
        \begin{itemize}
          \item Empfänger kann Paketfehler korrigieren ohne Sendewiederholung anzufordern
          \item Sender sendet hierzu redundante Pakete
        \end{itemize}
      \item Vorteil
        \begin{itemize}
          \item Reduziert die Verzögerung bis Daten beim Dienstnutzer abgeliefert werden können
        \end{itemize}
      \item Nachteil
        \begin{itemize}
          \item Redundante Information wird unabhängig von evtl. auftretendem Fehler gesendet
        \end{itemize}
      \item Hier: "verlorenes" Paket
        \begin{itemize}
          \item Paket kommt nie beim Empfänger an oder
          \item Paket kommt zu spät beim Empfänger an
        \end{itemize}
    \end{itemize}
